(1) 下端 \(z=0\) 张力为零，距下端 \(z\) 的截面需支撑下方重量，故 \(T(z)=\mu gz\)。小横向位移满足 \(\mu y_{tt}=\partial_z(Ty_z)=\mu g(z y_z)_z\)，即 \(y_{tt}=g(z y_{zz}+y_z)\)。张力随 \(z\) 变，不能写成恒速普通弦波方程。
(2) 分离后 \(zY''+Y'+(\omega^2/g)Y=0\)。置 \(\xi=2\omega\sqrt{z/g}\)，可化为 \(\xi^2Y_{\xi\xi}+\xi Y_\xi+\xi^2Y=0\)，其独立解为 \(J_0(\xi)\) 和 \(Y_0(\xi)\)。在 \(z\to0^+\) 时第二种 Neumann 函数对数发散，不满足有限位移与可接受端点正则性，故只保留 \(Y(z)=A J_0(2\omega\sqrt{z/g})\)。
(3) 固定顶端条件 \(Y(L)=0\) 使 \(J_0(2\omega\sqrt{L/g})=0\)。因此 \(\omega_n=(j_{0,n}/2)\sqrt{g/L}\)，\(n=1,2,\ldots\)；本征频率间距由 Bessel 零点决定而不是普通均匀张力弦的整数倍谱。
(4) 全系统的小振动能量 \(E=\frac12\int_0^L[\mu y_t^2+\mu gz\,y_z^2]dz\)。对时间微分并分部积分，边界功率 \([T y_z y_t]_0^L\)：顶端 \(y_t(L)=0\)，自由下端 \(T(0)=0\)，故 \(\dot E=0\)。下端不是固支端，不应取 \(y(0)=0\)；当 \(L\) 增大时所有频率按 \(L^{-1/2}\) 下降。
